\clearpage
\section{Consequences and the scale of the missing estimate}\label{sec:consequences}
\subsection{An asymptotic Fermat consequence}
The following implication is elementary and illustrates how repeated powers amplify the conjecture.
\begin{theorem}[Conditional large-exponent Fermat statement]
Assume Conjecture~\ref{conj:abc}. There is an integer $n_0$ such that, for every integer $n>n_0$, the equation
\[
 x^n+y^n=z^n
\]
has no positive integer solution.
\end{theorem}
\begin{proof}
Divide a putative solution by its common gcd. The resulting $x,y,z$ are pairwise coprime. Apply $abc$ to $a=x^n$, $b=y^n$, $c=z^n$ with the fixed exponent $\eps=1$. Then
\[
 z^n\le K_1\rad(xyz)^2\le K_1(xyz)^2<K_1z^6,
\]
because $x,y<z$. As $z\ge2$, this implies $2^{n-6}<K_1$ whenever $n>6$. Any fixed integer larger than $6+\log_2\max(1,K_1)$ excludes all such exponents.
\end{proof}
The theorem does not give Fermat's Last Theorem for every $n\ge3$. It supplies only a sufficiently-large-exponent implication, with a threshold depending on a conjectural constant. Fermat's Last Theorem is already an unconditional theorem. Its proof and the subsequent modularity theorem for elliptic curves form established arithmetic context~\cite{Wiles,TaylorWiles,BCDT}.

Using a smaller fixed $\eps$ sharpens the displayed exponent six toward three. It still leaves a constant-dependent threshold. Letting $\eps$ vary with $n$ without controlling $K_\eps$ would not be a valid way to remove the remaining finite exponents.

\subsection{Elliptic curves and the modified Szpiro comparison}
The marked Frey curve
\[
 E:y^2=x(x-a)(x+b)
\]
translates the additive triple into an elliptic curve. Its discriminant sees repeated depth, while its conductor sees prime support, with local model corrections at the small primes. This makes elliptic height estimates a natural companion to $abc$.

The precise equivalent statement uses more than the discriminant alone. For an elliptic curve over $\Q$, let $N$ be its conductor and let $c_4,c_6$ be the invariants of a global minimal integral Weierstrass equation. The modified Szpiro assertion is
\[
 \max\{|c_4|^3,c_6^2\}\le C_\eps N^{6+\eps}
\]
for every $\eps>0$, with a constant uniform in the elliptic curve. Oesterl\'e identifies this form as equivalent to the full $abc$ conjecture~\cite{Oesterle1988}.

This equivalence is cited as an external theorem. We do not derive its full small-prime and minimal-model accounting here. The ordinary statement involving only $|\Delta_{\min}|$ should not be silently substituted for the max-invariant formulation. A proposed project route through the Frey curve must preserve the marked integral model and the normalization of its differential, rather than only its rational isogeny class.

\subsection{Why local algebra is substantial but insufficient}
Several exact statements can coexist while the final height bound remains open:
\begin{enumerate}
\item A prime-power contact module can be defined canonically.
\item Its local carry structure can be computed exactly.
\item A Kummer class can carry nearly all of its finite order globally.
\item A logarithmic character can record the residual coherence.
\item A formal or analytic determinant can have large contact and be nonzero.
\end{enumerate}
None of these statements bounds the archimedean size, denominators, or arithmetic degree of a single realization that retains all the contact. The point where those costs are compared with $R_aR_b$ is the point at which the conjecture's quantitative content enters.

This distinction motivates the research's structural limitations. They eliminate specific ways of overlooking a cost: raw coordinate splitting, a split associated graded source, the full Fermat polarization, a complex period without an integral lattice, or a rational truncation without output contact. Each limitation narrows the requirements of the next valid construction.

\subsection{The research question to carry forward}
The direct output-sector target is
\[
 M_c\le K_\eps R_aR_bR^\eps.
\]
It asks for a uniform explanation of output repetition in terms of input support. The coherence route factors $M_c$ exactly and asks for a bound on the residual product $WC$. The determinant route seeks a nonzero arithmetic witness whose divisibility and global height imply the same inequality.

These are explicit interfaces, not estimates already established. A new local formula is progress within its own scope. A new global gain must reduce the exponent or comparison cost on the correctly normalized defect $c/R$.
